\begin{lemma}
\label{editorial-source-lemma-annihilator-flat-base-change}
Let $R \to S$ be a flat ring map. Let $M$ be an $R$-module and
$m \in M$. Then $\text{Ann}_R(m) S = \text{Ann}_S(m \otimes 1)$.
If $M$ is a finite $R$-module, then
$\text{Ann}_R(M)S=\text{Ann}_S(M\otimes_R S)$.
The latter conclusion also holds if finitely many elements
$m_1,\ldots,m_n\in M$ have
$\operatorname{Ann}_R(M)=\bigcap_{i=1}^n\operatorname{Ann}_R(m_i)$,
without requiring that they generate $M$.
\end{lemma}

\begin{proof}
Write the original ring map as $\varphi:R\to S$ and retain
$I=\operatorname{Ann}_R(m)$. The exact sequence
$$
0\longrightarrow I\longrightarrow R\longrightarrow M,
\qquad r\longmapsto rm,
$$
remains exact after tensoring with the flat $R$-module $S$.
Under $R\otimes_RS\to S$, $r\otimes s\mapsto\varphi(r)s$,
the last map becomes $s\mapsto m\otimes s=(m\otimes1)s$.
The image of $I\otimes_RS\to S$, $a\otimes s\mapsto\varphi(a)s$,
is exactly the extended ideal $IS$. Its equality with the
kernel therefore proves
$$
\operatorname{Ann}_R(m)S=\operatorname{Ann}_S(m\otimes1).
$$

For the original finite generating family $m_1,\ldots,m_n$,
retain $I_i=\operatorname{Ann}_R(m_i)$. An element kills $M$
exactly when it kills every generator. The elements
$m_i\otimes1$ generate $M\otimes_RS$, so the full comparison is
\[

\begin{split}
\operatorname{Ann}_S(M\otimes_RS)
&=\bigcap_{i=1}^n\operatorname{Ann}_S(m_i\otimes1)\\
&=\bigcap_{i=1}^n(I_iS)
=\left(\bigcap_{i=1}^n I_i\right)S
=\operatorname{Ann}_R(M)S.
\end{split}

\]
The finite-intersection equality is
Lemma \ref{editorial-source-lemma-flat-intersect-ideals}.

For the stronger assertion, suppose only that the original
finite list detects the annihilator:
$$
J=\operatorname{Ann}_R(M)=\bigcap_{i=1}^n I_i.
$$
These elements need not generate $M$. The sequence
$$
0\longrightarrow J\longrightarrow R\longrightarrow M^n,
\qquad r\longmapsto(rm_1,\ldots,rm_n),
$$
is exact. Tensoring and using the coordinate maps
$M^n\otimes_RS\simeq(M\otimes_RS)^n$ gives
$JS=\bigcap_i\operatorname{Ann}_S(m_i\otimes1)$.
Consequently every element annihilating the whole tensor
module belongs to $JS$. Conversely each finite sum
$\sum_j\varphi(a_j)s_j$ with $a_j\in J$ kills every tensor:
$$
\left(\sum_j\varphi(a_j)s_j\right)(m\otimes s)
=\sum_j(a_jm)\otimes s_js=0.
$$
This proves equality for the weaker hypothesis. For $n=0$
the empty intersection is $R$, forcing $M=0$; both module
annihilators are then the whole rings, as required.
For example, any nonempty free module, even with an infinite
basis, has its annihilator detected by one original basis
vector: $re=0$ implies $r=0$ by that coordinate.

\medskip\noindent
For an arbitrary module the same element calculation gives
the exact formula
$$
\operatorname{Ann}_S(M\otimes_RS)
=\bigcap_{m\in M}\bigl(\operatorname{Ann}_R(m)S\bigr),
$$
since all the elements $m\otimes1$ generate the tensor module.
Extension of ideals need not commute with this infinite
intersection, even for a faithfully flat ring map.

Retain the explicit example
$$
R=\mathbf Z,\qquad
M=\bigoplus_{n\geq1}\mathbf Z/2^n\mathbf Z,\qquad
S=\mathbf Z\times\mathbf Q,
\qquad \varphi(a)=(a,a).
$$
The rational numbers are the localization at the nonzero
integers. For every module $N$ the tensor-localization map is
$N\otimes_{\mathbf Z}\mathbf Q\to
(\mathbf Z\setminus\{0\})^{-1}N$,
$n\otimes a/b\mapsto an/b$, with inverse $n/b\mapsto n\otimes1/b$.
Localization preserves injections: a fraction mapping to zero
has a nonzero integer witness killing its numerator, and
injectivity gives the same relation in the original module.
Thus $\mathbf Q$ is flat, as also established in
Example \ref{editorial-source-example-flat-infinite-intersections}.
The exact tensor comparison
$$
N\otimes_{\mathbf Z}(\mathbf Z\times\mathbf Q)
\longrightarrow N\oplus(N\otimes_{\mathbf Z}\mathbf Q),
\qquad n\otimes(a,q)\longmapsto(an,n\otimes q)
$$
has inverse
$$
\left(n,\sum_j n_j\otimes q_j\right)
\longmapsto n\otimes(1,0)+\sum_j n_j\otimes(0,q_j).
$$
Both are balanced and their composites are identities on
the displayed generators. This proves flatness of $S$ and
faithfulness: the tensor of a homomorphism $h$ has first
component $h$, so it is zero only if $h=0$.

The original annihilator is
$$
\operatorname{Ann}_{\mathbf Z}(M)
=\bigcap_{n\geq1}2^n\mathbf Z=0,
$$
because no nonzero integer is divisible by every $2^n$.
Each original $m\in M$ has finite support, hence is killed
by some $2^N$. Therefore
$m\otimes q=(2^Nm)\otimes(q/2^N)=0$ in
$M\otimes_{\mathbf Z}\mathbf Q$. The preceding comparison gives
$M\otimes_{\mathbf Z}S\simeq M\oplus0$, with the action
$(a,q)(m,0)=(am,0)$. It follows exactly that
$$
\operatorname{Ann}_{\mathbf Z}(M)S=0,
\qquad
\operatorname{Ann}_S(M\otimes_{\mathbf Z}S)=0\times\mathbf Q\ne0.
$$
Every finite family in $M$ is killed by a common nonzero
power of $2$, so no such family detects its zero annihilator.
This same module has support $\{(2)\}$: its first summand
has nonzero localization at $(2)$, while at every other
prime $2$ is a unit and every element is killed by a power
of $2$. Thus its support is closed but is strictly smaller
than $V(\operatorname{Ann}_{\mathbf Z}(M))=\Spec(\mathbf Z)$.

The counterexample even has a finitely presented faithfully
flat algebra version. Put $S'=\mathbf Z\times\mathbf Z[1/2]$.
The identical localization and direct-sum arguments prove
faithful flatness and give annihilator $0\times\mathbf Z[1/2]$.
More explicitly,
$$
\mathbf Z[T]/(2T^2-T)\longrightarrow S',
\qquad f(T)\longmapsto(f(0),f(1/2)),
$$
is an isomorphism. The ideals $(T)$ and $(2T-1)$ are
comaximal since $2T-(2T-1)=1$; hence their intersection
is their product $(2T^2-T)$. The Chinese remainder map
is surjective with this kernel. Its factors are $\mathbf Z$
and $\mathbf Z[1/2]$: in the second quotient $2T=1$, and
the inverse to evaluation at $1/2$ sends $a/2^j$ to $aT^j$.
This proves the asserted presentation with the original
integer factors retained.

Finally, faithful flatness does always imply the contraction
identity
$$
\varphi^{-1}\bigl(\operatorname{Ann}_S(M\otimes_RS)\bigr)
=\operatorname{Ann}_R(M).
$$
Multiplication by $r$ on $M$ becomes multiplication by
$\varphi(r)$ on its tensor module, and a faithfully flat
tensor functor detects whether this homomorphism is zero.
Thus the counterexample fails extension, while contraction
still recovers the original annihilator. For $M=0$ or the
zero ring, all the annihilator statements retain these meanings.
\end{proof}